\subsection{CPython Allocation Optimizations}

A tuple is structurally simpler than a list. It has a fixed length, so CPython does not need to maintain spare capacity for future appends. Once the tuple exists, its slot count is fixed.

This gives tuples two practical implementation advantages:

\begin{itemize}
    \item the tuple object can be allocated with exactly the needed number of reference slots;
    \item CPython can apply allocation shortcuts for some tuple cases.
\end{itemize}

These are CPython implementation details, not general Python-language promises. Other Python implementations may use different internal layouts.

\begin{verbatim}
import platform
import sys

data = ("spam", "eggs", 42)

print(f"implementation: {platform.python_implementation()}")
print(f"python version:  {sys.version.split()[0]}")

print()
print(f"data: {data!r}")
print(f"type(data): {type(data)}")
print(f"len(data): {len(data)}")
print(f"size: {sys.getsizeof(data)} bytes")
\end{verbatim}

Possible output on one 64-bit CPython build:

\begin{verbatim}
implementation: CPython
python version:  3.12.3

data: ('spam', 'eggs', 42)
type(data): <class 'tuple'>
len(data): 3
size: 64 bytes
\end{verbatim}

The exact byte count can differ between Python versions, builds, and platforms. The important structural point is that a tuple does not carry list-style spare capacity for later growth.

\subsubsection{CPython Allocation Caching: Version-Dependent Recycling Optimizations for Small Tuple Objects}

CPython has historically used special allocation behavior for tuples. One simple case is the empty tuple. CPython reuses the empty tuple object.

\begin{verbatim}
a = ()
b = ()
c = tuple()

print(f"a: {a!r}")
print(f"b: {b!r}")
print(f"c: {c!r}")

print()
print(f"a is b: {a is b}")
print(f"a is c: {a is c}")

print()
print(f"id(a): {id(a)}")
print(f"id(b): {id(b)}")
print(f"id(c): {id(c)}")
\end{verbatim}

Possible output on CPython:

\begin{verbatim}
a: ()
b: ()
c: ()

a is b: True
a is c: True

id(a): 140619796135544
id(b): 140619796135544
id(c): 140619796135544
\end{verbatim}

This is safe because the empty tuple cannot be changed. There is no internal slot that could later be reassigned.

CPython also performs recycling optimizations for some non-empty tuple allocations. The exact rules are implementation details and may change between versions. Therefore, tuple identity reuse must never be used as program logic.

The following demonstration may show object-id reuse on CPython, but it is intentionally presented as an observation experiment, not as a guarantee.

\begin{verbatim}
import gc

ids = []

for nr in range(10):
    data = (nr, nr + 1)
    ids.append(id(data))
    del data
    gc.collect()

print("observed ids")
for value in ids:
    print(value)

print()
print(f"unique id count: {len(set(ids))}")
print(f"total creations: {len(ids)}")
\end{verbatim}

Possible output on one CPython run:

\begin{verbatim}
observed ids
140619786871232
140619786871232
140619786871232
140619786871232
140619786871232
140619786871232
140619786871232
140619786871232
140619786871232
140619786871232

unique id count: 1
total creations: 10
\end{verbatim}

Another run, another Python version, or another implementation may show different results. The only valid lesson is that CPython may recycle tuple memory after a tuple is destroyed.

The same warning applies to identity comparisons.

\begin{verbatim}
a = (10, 20)
b = (10, 20)

print(f"a == b: {a == b}")
print(f"a is b: {a is b}")

print()
print(f"id(a): {id(a)}")
print(f"id(b): {id(b)}")
\end{verbatim}

Possible output:

\begin{verbatim}
a == b: True
a is b: False

id(a): 140619786871232
id(b): 140619786871488
\end{verbatim}

The tuples have equal contents. They should be compared with \texttt{==}, not \texttt{is}.

Some constant tuples may appear to be shared because the compiler can store constants and reuse them inside code objects. This is another implementation detail.

\begin{verbatim}
a = (1, 2, 3)
b = (1, 2, 3)

print(f"a == b: {a == b}")
print(f"a is b: {a is b}")

print()
print(f"id(a): {id(a)}")
print(f"id(b): {id(b)}")
\end{verbatim}

Possible output on one CPython run:

\begin{verbatim}
a == b: True
a is b: True

id(a): 140619786873024
id(b): 140619786873024
\end{verbatim}

This result is not the tuple equality rule. It is a storage optimization effect. The semantic rule is still:

\begin{quote}
Use \texttt{==} to compare tuple contents. Do not use \texttt{is} to test whether two tuples contain the same values.
\end{quote}

Tuple construction from an existing exact tuple may also return the same object, because there is no reason to copy an immutable tuple.

\begin{verbatim}
a = ("Homer", "Marge")
b = tuple(a)

print(f"a: {a!r}")
print(f"b: {b!r}")

print()
print(f"a == b: {a == b}")
print(f"a is b: {a is b}")
\end{verbatim}

Possible output:

\begin{verbatim}
a: ('Homer', 'Marge')
b: ('Homer', 'Marge')

a == b: True
a is b: True
\end{verbatim}

This is safe because neither \texttt{a} nor \texttt{b} can mutate the tuple structure.

The practical rule is:

\begin{quote}
Tuple allocation optimizations are allowed because tuples are structurally immutable. They are useful for CPython performance, but they must not be treated as part of normal program meaning.
\end{quote}

\subsubsection{Memory Footprint Metrics: Comparing Fixed \texttt{tuple} Layouts Against the Dynamic Tracking Buffers of \texttt{list}}

A list must support growth. For that reason, CPython lists maintain dynamic internal storage and may reserve extra pointer slots beyond the visible length.

A tuple does not support growth. Its fixed slot count can be allocated directly.

The difference can be measured with \texttt{sys.getsizeof()}.

\begin{verbatim}
import sys

values = [10, 20, 30, 42]

items_list = [10, 20, 30, 42]
items_tuple = (10, 20, 30, 42)

print(f"items_list:  {items_list!r}")
print(f"items_tuple: {items_tuple!r}")

print()
print(f"len(items_list):  {len(items_list)}")
print(f"len(items_tuple): {len(items_tuple)}")

print()
print(f"sys.getsizeof(items_list):  {sys.getsizeof(items_list)}")
print(f"sys.getsizeof(items_tuple): {sys.getsizeof(items_tuple)}")
\end{verbatim}

Possible output on one 64-bit CPython build:

\begin{verbatim}
items_list:  [10, 20, 30, 42]
items_tuple: (10, 20, 30, 42)

len(items_list):  4
len(items_tuple): 4

sys.getsizeof(items_list):  88
sys.getsizeof(items_tuple): 72
\end{verbatim}

The exact numbers are not universal. The usual pattern is that a tuple with the same number of elements is smaller than a list because the tuple does not need dynamic spare capacity.

A size table makes the pattern clearer.

\begin{verbatim}
import sys

print(f"{'n':>2} {'list':>8} {'tuple':>8}")

for n in range(0, 11):
    items_list = list(range(n))
    items_tuple = tuple(range(n))

    list_size = sys.getsizeof(items_list)
    tuple_size = sys.getsizeof(items_tuple)

    print(f"{n:>2} {list_size:>8} {tuple_size:>8}")
\end{verbatim}

Possible output on one 64-bit CPython build:

\begin{verbatim}
 n     list    tuple
 0       56       40
 1       72       48
 2       72       56
 3       88       64
 4       88       72
 5      104       80
 6      104       88
 7      120       96
 8      120      104
 9      136      112
10      136      120
\end{verbatim}

The tuple size usually grows in a direct slot-count pattern. The list size often grows in jumps because of allocation rounding and spare capacity.

This difference becomes more visible when a list is grown gradually with \texttt{append()}.

\begin{verbatim}
import sys

items = []

print(f"{'len':>3} {'list size':>10}")

last_size = sys.getsizeof(items)
print(f"{len(items):>3} {last_size:>10}")

for nr in range(1, 13):
    items.append(nr)
    size = sys.getsizeof(items)

    if size != last_size:
        print(f"{len(items):>3} {size:>10}")
        last_size = size
\end{verbatim}

Possible output on one 64-bit CPython build:

\begin{verbatim}
len  list size
  0         56
  1         88
  5        120
  9        184
\end{verbatim}

The list size does not increase at every length. It increases when the internal storage is resized.

A tuple of the same visible lengths does not need this spare-capacity behavior.

\begin{verbatim}
import sys

print(f"{'len':>3} {'tuple size':>10}")

for n in [0, 1, 5, 9]:
    data = tuple(range(n))
    print(f"{n:>3} {sys.getsizeof(data):>10}")
\end{verbatim}

Possible output on one 64-bit CPython build:

\begin{verbatim}
len tuple size
  0         40
  1         48
  5         80
  9        112
\end{verbatim}

This does not mean tuples are always the correct replacement for lists. The choice depends on the required behavior.

\begin{itemize}
    \item Use a list when the sequence must be edited in place.
    \item Use a tuple when the sequence represents a fixed structure.
\end{itemize}

The memory comparison is only one consequence of the deeper structural difference.

The following example shows the behavioral reason first.

\begin{verbatim}
items_list = ["Jerry", "Elaine"]
items_tuple = ("Jerry", "Elaine")

items_list.append("George")

print(f"items_list: {items_list!r}")

try:
    items_tuple.append("George")
except AttributeError as err:
    print("tuple append failed")
    print(f"message: {err}")
\end{verbatim}

Possible output:

\begin{verbatim}
items_list: ['Jerry', 'Elaine', 'George']
tuple append failed
message: 'tuple' object has no attribute 'append'
\end{verbatim}

The list pays for editability with dynamic storage machinery. The tuple avoids that machinery because its structure is fixed.

The practical rule is:

\begin{quote}
A tuple is usually smaller because it stores a fixed number of reference slots. A list is more flexible because it owns a dynamically resized slot buffer. Choose between them by meaning first: editable sequence or fixed record-like sequence.
\end{quote}